\section*{Hoofdstuk \ref{ch:b1:topology} --- Topologie van de reële rechte}

\begin{solution}{exo:b1:topology:1}
$\intoo{0}{1} \cup \intoo{2}{3}$: \omterm{def:b1:topology:open}{open} (vereniging van \omterm{def:b1:topology:open}{open
verzamelingen}), niet \omterm{def:b1:topology:closed}{gesloten} ($\frac 1n \to 0$ ligt erbuiten).

$\intco{0}{1}$: geen van beide. Niet \omterm{def:b1:topology:open}{open} (geen \omterm{prop:b1:reals:intervals}{interval} rond $0$ ligt
erbinnen); niet \omterm{def:b1:topology:closed}{gesloten} ($1 - \frac1n \to 1 \notin$ de \omterm{def:b1:logic:sets}{verzameling}).

$\{0\} \cup \intcc{1}{2}$: \omterm{def:b1:topology:closed}{gesloten} (eindige vereniging van \omterm{def:b1:topology:closed}{gesloten
verzamelingen}), niet \omterm{def:b1:topology:open}{open} (faalt in $0$).

$\R \setminus \Z$: \omterm{def:b1:topology:open}{open} ($\Z$ is \omterm{def:b1:topology:closed}{gesloten}), niet \omterm{def:b1:topology:closed}{gesloten}: de rij
$\bigl(\frac 1n\bigr)$ ligt erin, maar haar limiet $0$ behoort tot $\Z$
en ontsnapt dus aan de \omterm{def:b1:logic:sets}{verzameling}.

$\Q \cap \intoo{0}{1}$: geen van beide. Niet \omterm{def:b1:topology:open}{open}: elk \omterm{prop:b1:reals:intervals}{interval} rond een
rationaal getal bevat irrationale getallen. Niet \omterm{def:b1:topology:closed}{gesloten}: ze bevat rijen
die naar het irrationale $\frac{\sqrt 2}{2}$ gaan (\omterm{def:b1:topology:dense}{dichtheid}).
\end{solution}

\begin{solution}{exo:b1:topology:2}
$A = \intoc{0}{1} \cup \{2\}$: $\mathring A = \intoo{0}{1}$, $\overline A
= \intcc{0}{1} \cup \{2\}$, $\partial A = \{0, 1, 2\}$.

$A = \R \setminus \Q$: $\mathring A = \emptyset$ (elk \omterm{prop:b1:reals:intervals}{interval} bevat
rationale getallen), $\overline A = \R$ (\omterm{def:b1:topology:dense}{dichtheid} van de irrationale
getallen), $\partial A = \R$.

$A = \bigl\{\frac{(-1)^n n}{n+1}\bigr\}$: de termen met even index gaan
naar $1$, die met oneven index naar $-1$, en geen van beide behoort tot
$A$. $\mathring A = \emptyset$ (geïsoleerde punten), $\overline A = A
\cup \{-1, 1\}$, $\partial A = \overline A$.
\end{solution}

\begin{solution}{exo:b1:topology:3}
\emph{Complementen.} $F = \{a_1 < a_2 < \dots < a_k\}$: het complement is
de vereniging van de \omterm{def:b1:topology:open}{open} \omterm{prop:b1:reals:intervals}{intervallen} $\intoo{-\infty}{a_1}$,
$\intoo{a_i}{a_{i+1}}$ en $\intoo{a_k}{+\infty}$ --- \omterm{def:b1:topology:open}{open} volgens
\cref{prop:b1:topology:openstable}.

\emph{Rijen.} Zij $u_n \in F$ met $u_n \to \ell$. Met $\varepsilon =
\min\{\abs{a_i - a_j} : i \neq j\}/2 > 0$ (of een willekeurige
$\varepsilon$ als $F$ een singleton is): voorbij een zekere rang liggen
alle termen binnen $\varepsilon$ van $\ell$ en dus binnen $2\varepsilon$
van elkaar, wat afdwingt dat ze vanaf die rang één enkele $a_i$ zijn; dan
is $\ell = a_i \in F$.
\end{solution}

\begin{solution}{exo:b1:topology:4}
$\subseteq$: $\overline A \cup \overline B$ is \omterm{def:b1:topology:closed}{gesloten} (eindige
vereniging) en bevat $A \cup B$, dus bevat ze de \emph{kleinste} \omterm{def:b1:topology:closed}{gesloten}
oververzameling $\overline{A \cup B}$. $\supseteq$: uit $A \subseteq A
\cup B$ volgt $\overline A \subseteq \overline{A \cup B}$ (de \omterm{def:b1:topology:closure}{afsluiting}
is monotoon: adherente punten van $A$ zijn adherent aan de grotere
\omterm{def:b1:logic:sets}{verzameling}), en analoog voor $B$.

Tegenvoorbeeld voor doorsneden: $A = \intoo{0}{1}$ en $B =
\intoo{1}{2}$: $\overline{A \cap B} = \overline\emptyset = \emptyset$
terwijl $\overline A \cap \overline B = \{1\}$.
\end{solution}

\begin{solution}{exo:b1:topology:5}
Gebruik de karakterisering met rijen
(\cref{thm:b1:topology:seqclosed}). Zij $S = \{u_n\} \cup \{\ell\}$ en
$(v_k)$ een rij in $S$ met $v_k \to m$; toon $m \in S$ aan. Twee
gevallen. Wordt een zekere waarde $v \in S$ oneindig vaak door $(v_k)$
aangenomen, dan geeft een constante \omterm{def:b1:seq:subsequence}{deelrij} $m = v \in S$. Anders wordt
elke waarde eindig vaak aangenomen; in het bijzonder verschijnt voor elke
$n$ de term $u_n$ eindig vaak, en $\ell$ eveneens. Dan zijn er voor elke
$N$ maar eindig veel indices $k$ met $v_k \in \{u_0, \dots, u_N, \ell\}$:
de overige $v_k$ zijn termen $u_n$ met $n > N$. Kies bij gegeven
$\varepsilon > 0$ een $N$ met $\abs{u_n - \ell} \leq \varepsilon$ voor $n
> N$: op eindig veel na voldoen alle $v_k$ aan $\abs{v_k - \ell} \leq
\varepsilon$. Dus $v_k \to \ell$, waaruit $m = \ell \in S$.
\end{solution}

\begin{solution}{exo:b1:topology:6}
$U + A = \bigcup_{a \in A} (U + a)$, en elke verschuiving $U + a$ is \omterm{def:b1:topology:open}{open}
(verschuif de intervalbewijsstukken). Een vereniging van \omterm{def:b1:topology:open}{open
verzamelingen} is \omterm{def:b1:topology:open}{open} (\cref{prop:b1:topology:openstable}).

\omterm{def:b1:topology:closed}{Gesloten verzamelingen}: $\Z$ en $\sqrt 2\,\Z$ zijn \omterm{def:b1:topology:closed}{gesloten} (als
$\alpha\Z$: convergente rijen zijn uiteindelijk constant, vergelijk
\cref{ex:b1:topology:closed}). Hun som $\Z + \sqrt 2\,\Z$ ligt \omterm{def:b1:topology:dense}{dicht} in
$\R$ (\cref{exo:b1:reals:9}) maar is niet $\R$ (ze is aftelbaar, of
eenvoudiger: $\frac{\sqrt 2}{2} \notin \Z + \sqrt2\Z$, want anders zou
$\sqrt 2$ rationaal zijn --- uit $\frac{\sqrt2}{2} = m + n\sqrt 2$ volgt
$(2n - 1)\sqrt 2 = -2m$, dus $\sqrt 2 \in \Q$ tenzij $n = \frac12$, wat
onmogelijk is). Een \omterm{def:b1:topology:dense}{dichte} echte deelverzameling is niet \omterm{def:b1:topology:closed}{gesloten}: haar
\omterm{def:b1:topology:closure}{afsluiting} is $\R$ en dus niet zijzelf.
\end{solution}

\begin{solution}{exo:b1:topology:7}
$\Z$: de \omterm{def:b1:topology:open}{omgeving} $\intoo{n - \frac12}{n + \frac12}$ van $n$ ontmoet $\Z$
alleen in $n$.

$A = \{\frac1n : n \in \N^*\}$: rond $\frac 1n$ isoleert het \omterm{prop:b1:reals:intervals}{interval} met
straal $\frac{1}{n} - \frac{1}{n+1} = \frac{1}{n(n+1)}$ (bijvoorbeeld
gehalveerd) het punt van zijn buren --- alle punten zijn geïsoleerd. En
toch is $0 \in \overline A \setminus A$: geïsoleerde punten beletten
adherente buitenstaanders niet.

Zijn alle punten van $A$ geïsoleerd, dan is geen enkel punt van $A$
\omterm{def:b1:topology:closure}{inwendig}, want een \omterm{def:b1:topology:closure}{inwendig} punt heeft een heel \omterm{prop:b1:reals:intervals}{interval} $A$-buren om
zich heen (een \omterm{prop:b1:reals:intervals}{interval} is oneindig), in tegenspraak met de isolatie. Dus
is $\mathring A = \emptyset$.
\end{solution}

\begin{solution}{exo:b1:topology:8}
Is $A \subseteq \intcc{-M}{M}$, dan bevat de \omterm{def:b1:topology:closed}{gesloten verzameling}
$\intcc{-M}{M}$ de \omterm{def:b1:logic:sets}{verzameling} $A$ en dus ook $\overline A$ (de kleinste
\omterm{def:b1:topology:closed}{gesloten} oververzameling): $\overline A$ is begrensd.

Zij $s = \sup A$ (eindig). Omdat $A \subseteq \overline A$, is $\sup
\overline A \geq s$. Omgekeerd is $\overline A \subseteq
\intoc{-\infty}{s}$: die halfrechte is \omterm{def:b1:topology:closed}{gesloten} en bevat $A$; dus is elk
element van $\overline A$ hoogstens $s$, waaruit $\sup\overline A \leq
s$. Gelijkheid.

$s \in \overline A$: volgens de $\varepsilon$-karakterisering
(\cref{prop:b1:reals:epsilon}) bevat elk \omterm{prop:b1:reals:intervals}{interval} $\intoo{s -
\varepsilon}{s + \varepsilon}$ een element van $A$: $s$ is adherent.
\end{solution}

\begin{solution}{exo:b1:topology:9}
Stel dat $A$ \omterm{def:b1:topology:open}{open} en \omterm{def:b1:topology:closed}{gesloten} is, met $a \in A$ en $b \in \R \setminus
A$; zonder verlies van algemeenheid $a < b$. De \omterm{def:b1:logic:sets}{verzameling} $B = A \cap
\intcc{a}{b}$ is niet-leeg ($a$) en begrensd: zij $s = \sup B$. Volgens
\cref{exo:b1:topology:8} is $s \in \overline B \subseteq \overline A = A$
($A$ \omterm{def:b1:topology:closed}{gesloten}). Merk op dat $s \leq b$, en omdat $b \notin A$ is zelfs $s
< b$. Nu is $A$ \omterm{def:b1:topology:open}{open}: een zeker \omterm{prop:b1:reals:intervals}{interval} $\intoo{s - r}{s + r}$ ligt in
$A$, en we mogen $r < b - s$ nemen. Dan behoort $s + \frac{r}{2}$ tot $A
\cap \intcc{a}{b} = B$ en overtreft het $s$ --- in tegenspraak met $s =
\sup B$. Zo'n paar $(a, b)$ bestaat dus niet: een van $A$ en $\R
\setminus A$ is leeg.
\end{solution}

\begin{solution}{exo:b1:topology:10}
$I_x$ is een vereniging van \omterm{def:b1:topology:open}{open} \omterm{prop:b1:reals:intervals}{intervallen} die alle $x$ bevatten: ze is
\omterm{def:b1:topology:open}{open}, en ze is een \omterm{prop:b1:reals:intervals}{interval} omdat ze convex is --- is $u < z < v$ met $u,
v \in I_x$, dan liggen $u$ en $v$ in \omterm{def:b1:topology:open}{open} deelintervallen $J_u \ni x$ en
$J_v \ni x$ van $U$, en $J_u \cup J_v$ is een \omterm{prop:b1:reals:intervals}{interval} (beide bevatten
$x$) binnen $U$ dat $z$ bevat; dus $z \in I_x$
(\cref{prop:b1:reals:intervals}).

Is $I_x \cap I_y \neq \emptyset$, dan is $I_x \cup I_y$ een \omterm{def:b1:topology:open}{open} \omterm{prop:b1:reals:intervals}{interval}
(convex: twee overlappende \omterm{prop:b1:reals:intervals}{intervallen}) bevat in $U$ dat $x$ en $y$
bevat, dus $I_x \cup I_y \subseteq I_x$ en $\subseteq I_y$ wegens de
maximaliteit van elk: $I_x = I_y$.

Dus is $U$ de disjuncte vereniging van de verschillende \omterm{def:b1:logic:sets}{verzamelingen}
$I_x$ (elke $x \in U$ ligt in haar eigen $I_x$). Aftelbaarheid: elk
niet-leeg \omterm{def:b1:topology:open}{open} \omterm{prop:b1:reals:intervals}{interval} $I$ van de familie bevat een rationaal getal
$q_I$ (\cref{thm:b1:reals:density}), en verschillende disjuncte
\omterm{prop:b1:reals:intervals}{intervallen} krijgen verschillende rationale getallen: de familie
injecteert in $\Q$, dat aftelbaar is (het wordt door paren gehele
getallen geïndexeerd). Er zijn dus hoogstens aftelbaar veel \omterm{prop:b1:reals:intervals}{intervallen}.
\end{solution}

\begin{solution}{exo:b1:topology:11}
$\overline A = A \cup A'$. ($\supseteq$) Steeds is $A \subseteq \overline
A$; en is $x \in A'$, dan ontmoet elke \omterm{def:b1:topology:open}{omgeving} van $x$ de \omterm{def:b1:logic:sets}{verzameling} $A
\setminus \{x\} \subseteq A$, dus is $x$ adherent. ($\subseteq$) Zij $x
\in \overline A$. Is $x \in A$, dan zijn we klaar. Is $x \notin A$, dan
ontmoet elke \omterm{def:b1:topology:open}{omgeving} van $x$ de \omterm{def:b1:logic:sets}{verzameling} $A = A \setminus \{x\}$: $x
\in A'$.

Bijgevolg is $A$ \omterm{def:b1:topology:closed}{gesloten} $\iff$ $A = \overline A = A \cup A'$ $\iff$ $A'
\subseteq A$.

$A = \{\frac 1n\}$: $0$ is een ophopingspunt ($\frac 1n \to 0$, met
termen $\neq 0$); elke $\frac 1n$ is geïsoleerd
(\cref{exo:b1:topology:7}) en dus niet in $A'$; en een punt $x \notin A
\cup \{0\}$ heeft een heel \omterm{prop:b1:reals:intervals}{interval} dat $A$ mijdt (tussen zijn twee buren
in $A \cup \{0\}$, of voorbij $1$). Dus $A' = \{0\}$.

$\Z' = \emptyset$: elk geheel getal is geïsoleerd, en elk niet-geheel
getal heeft een \omterm{def:b1:topology:open}{omgeving} binnen $\R \setminus \Z$.

$\Q' = \R$: elk \omterm{prop:b1:reals:intervals}{interval} rond een willekeurig reëel getal bevat oneindig
veel rationale getallen (\cref{thm:b1:reals:density}), in het bijzonder
één dat van het middelpunt verschilt.
\end{solution}

\begin{solution}{exo:b1:topology:12}
\begin{enumerate}
  \item Voor elke $a \in A$ is $\abs{x - a} \leq \abs{x - y} + \abs{y -
        a}$, dus $d_A(x) \leq \abs{x - y} + \abs{y - a}$; het \omterm{def:b1:reals:bounds}{infimum}
        over $a$ nemen geeft $d_A(x) \leq \abs{x - y} + d_A(y)$. $x$ en
        $y$ verwisselen geeft de andere ongelijkheid: $\abs{d_A(x) -
        d_A(y)} \leq \abs{x - y}$.
  \item $d_A(x) = 0$ $\iff$ er is voor elke $\varepsilon > 0$ een $a \in
        A$ met $\abs{x - a} < \varepsilon$ $\iff$ elk \omterm{prop:b1:reals:intervals}{interval} rond $x$
        ontmoet $A$ $\iff$ $x \in \overline A$. Is $F$ \omterm{def:b1:topology:closed}{gesloten} en $x
        \notin F = \overline F$, dan is $d_F(x) \neq 0$, dus $d_F(x) >
        0$.
  \item Is $d_F(x) < \varepsilon$, zet dan $r = \varepsilon - d_F(x) >
        0$: voor $\abs{y - x} < r$ geeft onderdeel (1) dat $d_F(y) \leq
        d_F(x) + \abs{x - y} < \varepsilon$: het \omterm{prop:b1:reals:intervals}{interval} $\intoo{x -
        r}{x + r}$ ligt in $V_\varepsilon$, die dus \omterm{def:b1:topology:open}{open} is; en ze bevat
        $F$ omdat $d_F$ daar nul is. Ten slotte is $x \in
        \bigcap_{\varepsilon>0} V_\varepsilon$ $\iff$ $d_F(x) <
        \varepsilon$ voor alle $\varepsilon$ $\iff$ $d_F(x) = 0$ $\iff$
        $x \in \overline F = F$. Omdat $\bigcap_{\varepsilon > 0}
        V_\varepsilon = \bigcap_{n \geq 1} V_{1/n}$, is elke \omterm{def:b1:topology:closed}{gesloten
        verzameling} een aftelbare doorsnede van \omterm{def:b1:topology:open}{open verzamelingen}.
\end{enumerate}
\end{solution}

\begin{solution}{pb:b1:topology:1}
\textbf{1.} $C_1 = \intcc{0}{\frac13} \cup \intcc{\frac23}{1}$ en
\[
C_2 = \intcc{0}{\tfrac19} \cup \intcc{\tfrac29}{\tfrac13}
\cup \intcc{\tfrac23}{\tfrac79} \cup \intcc{\tfrac89}{1} .
\]
Inductie: is $C_n$ een disjuncte vereniging van $2^n$ \omterm{def:b1:topology:closed}{gesloten} segmenten
van lengte $3^{-n}$, dan laat het verwijderen van het \omterm{def:b1:topology:open}{open} middelste
derde uit elk segment per ouder twee \omterm{def:b1:topology:closed}{gesloten} segmenten van lengte
$3^{-n-1}$ over: $2^{n+1}$ segmenten, twee aan twee disjunct (kinderen
van verschillende ouders zijn gescheiden omdat de ouders dat waren;
kinderen van één ouder zijn gescheiden door het verwijderde gat).

\textbf{2.} Elke $C_n$ is een eindige vereniging van segmenten en dus
\omterm{def:b1:topology:closed}{gesloten}; $C = \bigcap C_n$ is een doorsnede van \omterm{def:b1:topology:closed}{gesloten verzamelingen}:
\omterm{def:b1:topology:closed}{gesloten} (\cref{def:b1:topology:closed}); begrensd ($\subseteq
\intcc{0}{1}$): compact volgens \cref{thm:b1:topology:compact}.
Niet-leeg: $0$ ligt in het meest linkse segment van elke $C_n$. Zij $a$
een eindpunt van een segment $S$ van $C_n$. Voor $m \leq n$ is $a \in C_n
\subseteq C_m$. Voor de latere stadia: het verwijderen van het middelste
derde haalt nooit een eindpunt weg, en $a$ is opnieuw een eindpunt van
een van de twee kinderen van $S$ (dat aan $a$ raakt); met inductie is $a
\in C_m$ voor alle $m \geq n$: $a \in C$.

\textbf{3.} Totale lengte van $C_n$: $2^n \cdot 3^{-n} = (\frac23)^n \to
0$. Kies bij gegeven $\varepsilon > 0$ een $n$ met $(\frac23)^n \leq
\varepsilon$: dan is $C \subseteq C_n$, een vereniging van eindig veel
segmenten met totale lengte $\leq \varepsilon$.

\textbf{4.} Zij $I \subseteq C$ een \omterm{prop:b1:reals:intervals}{interval} met twee verschillende
punten. Voor elke $n$ is $I \subseteq C_n$, en $I$ moet, convex zijnde,
binnen één \emph{enkel} segment van $C_n$ liggen (twee segmenten
ontmoeten zou $I$ dwingen een punt van het gat ertussen te bevatten, dat
buiten $C_n$ ligt). De lengte van $I$ is dus $\leq 3^{-n}$ voor elke $n$:
tegenspraak. De enige \omterm{prop:b1:reals:intervals}{intervallen} binnen $C$ zijn dus leeg of singletons;
in het bijzonder past geen enkele $\intoo{x-r}{x+r}$ binnen $C$:
$\mathring C = \emptyset$. Omdat $C$ \omterm{def:b1:topology:closed}{gesloten} is, heeft $\overline C = C$
een leeg \omterm{def:b1:topology:closure}{inwendige}: $C$ is nergens \omterm{def:b1:topology:dense}{dicht}.

\textbf{5.} Schrijf $\varphi_0(x) = \frac x3$ en $\varphi_2(x) = \frac{2
+ x}{3}$, stijgende affiene bijecties van $\intcc{0}{1}$ op
$\intcc{0}{\frac13}$ en $\intcc{\frac23}{1}$. Bewering: $C_{n+1} =
\varphi_0(C_n) \cup \varphi_2(C_n)$. Voor $n = 0$ is dat vraag 1.
Inductie: een stijgende affiene \omterm{def:b1:logic:map}{afbeelding} stuurt het middelste derde van
een segment naar het middelste derde van het beeldsegment, zodat het
verwijderen van middelste derden met $\varphi_0$ en $\varphi_2$
verwisselt; de verwijderingsstap toepassen op $C_{n+1} = \varphi_0(C_n)
\cup \varphi_2(C_n)$ levert $C_{n+2} = \varphi_0(C_{n+1}) \cup
\varphi_2(C_{n+1})$. Doorsnijden over $n$: voor $x \leq \frac13$ is $x
\in C \iff x \in \varphi_0(C_n)$ voor alle $n$ $\iff 3x \in \bigcap C_n =
C$; analoog op $\intcc{\frac23}{1}$; en geen punt van
$\intoo{\frac13}{\frac23}$ ligt in $C_1$. Bijgevolg is $C = \varphi_0(C)
\cup \varphi_2(C)$, disjunct.

\textbf{6.} Inductie naar $n$; het geval $n = 0$ zegt dat elke $x \in
\intcc{0}{1}$ een code heeft, wat \cref{pb:b1:reals:1} is (vraag 9 voor
$x < 1$; en $1 = (0.\overline 2)_3$). Neem de equivalentie op rang $n$
aan. Is $x \in C_{n+1}$, dan is volgens vraag 5 $x = \varphi_i(z)$ met $z
\in C_n$ en $i \in \{0, 2\}$; is $(e_k)$ een code van $z$ waarvan de
eerste $n$ cijfers $1$-vrij zijn, dan heeft $(i, e_1, e_2, \dots)$ als
partiële sommen $\frac i3 + \frac13\sum_{k\leq m} e_k 3^{-k} \to
\varphi_i(z) = x$: een code van $x$ waarvan de eerste $n + 1$ cijfers
$1$-vrij zijn. Omgekeerd, heeft $x$ een code $(d_k)$ met $d_1, \dots,
d_{n+1} \in \{0, 2\}$, dan heeft de verschoven rij $(d_2, d_3, \dots)$
een zekere waarde $z \in \intcc{0}{1}$ waarvan de eerste $n$ cijfers
$1$-vrij zijn, en de berekening van de partiële sommen achterstevoren
gelezen geeft $x = \varphi_{d_1}(z)$; met de inductie is $z \in C_n$, dus
$x \in C_{n+1}$ wegens vraag 5. Ten slotte: een volledig $1$-vrije code
legt $x$ in elke $C_n$ en dus in $C$; omgekeerd, is $x \in C$, dan heeft
voor elke $n$ een van de hoogstens twee codes van $x$
(\cref{pb:b1:reals:1}, vraag 11) haar eerste $n$ cijfers $1$-vrij; één
vaste code moet dat voor willekeurig grote $n$ doen (duivenhok tussen
twee codes), en een code waarvan de eerste $n$ cijfers $1$-vrij zijn voor
willekeurig grote $n$ is ronduit $1$-vrij.

\textbf{7.} $0 = (0.\overline 0)_3$, $1 = (0.\overline 2)_3$, $\frac13 =
(0.0\overline{2})_3$ (de oneigenlijke tweeling van $(0.1)_3$) en
$\frac23 = (0.2\overline{0})_3$. Meetkundige sommen:
\[
(0.\overline{02})_3 = \sum_{j\geq1} \frac{2}{9^{\,j}}
= \frac{2/9}{1 - 1/9} = \frac14 ,
\qquad
(0.\overline{20})_3 = \sum_{j\geq1} \frac{2}{3\cdot 9^{\,j-1}}
= \frac{2/3}{1 - 1/9} = \frac34 ,
\]
beide $1$-vrij: $\frac14, \frac34 \in C$. (De oneindige sommen staan kort
voor \omterm{def:b1:reals:bounds}{suprema} van partiële sommen, als in \cref{pb:b1:reals:1}.) De
eindpunten van segmenten van $C_n$ hebben de vorm $m/3^n$ (inductie: de
eindpunten van kinderen zijn eindpunten van de ouder of verschillen er
een veelvoud van $3^{-n-1}$ van). Was $\frac14 = \frac{m}{3^n}$, dan zou
$3^n = 4m$, en $4 \nmid 3^n$: onmogelijk. Dus ligt $\frac14$ in $C$
zonder ooit een eindpunt te zijn.

\textbf{8.} Stel dat $x$ twee verschillende $1$-vrije codes had.
Überhaupt twee codes hebben betekent (\cref{pb:b1:reals:1}, vraag 11, in
grondtal $3$) dat $x = m/3^N \in \intoo{0}{1}$ en dat de twee codes de
afbrekende zijn, met laatste cijfer $d_N \in \{1, 2\}$ ongelijk aan nul
gevolgd door nullen, en haar tweeling, met $d_N - 1$ op positie $N$
gevolgd door tweeën. Is $d_N = 1$, dan bevat de eerste een $1$; is $d_N =
2$, dan draagt de tweeling $d_N - 1 = 1$. In beide gevallen is hoogstens
één van het paar $1$-vrij: tegenspraak. Elke $x \in C$ heeft dus precies
één $1$-vrije code (het bestaan volgt uit vraag 6), en verschillende
rijen over $\{0,2\}$ hebben verschillende waarden. Elke rij over
$\{0,2\}$ heeft een waarde in $\intcc{0}{1}$ (partiële sommen $\leq 1$)
waarvan alle beginstukken $1$-vrij zijn, dus een waarde in elke $C_n$,
dat wil zeggen in $C$: de waardeafbeelding is een bijectie van de rijen
over $\{0,2\}$ op $C$.

\textbf{9.} Zij $(d^{(k)})$ de $1$-vrije code van $x_k$ en zet $e_k = 2 -
d^{(k)}_k \in \{0, 2\}$: een rij over $\{0,2\}$ waarvan de waarde $y$ in
$C$ ligt met $(e_k)$ als unieke $1$-vrije code (vraag 8). Voor elke $k$
verschillen de codes van $y$ en $x_k$ op positie $k$, dus $y \neq x_k$:
geen enkele \omterm{def:b1:logic:map}{afbeelding} $\N^* \to C$ is \omterm{def:b1:logic:inj}{surjectief}. Een overaftelbare
\omterm{def:b1:logic:sets}{verzameling} van lengte nul: groot in \omterm{def:b1:counting:card}{kardinaliteit} en klein in maat,
tegelijk.

\textbf{10.} Compactheid: de geslotenheid van de oneindige doorsnede plus
de begrensdheid (vraag 2) --- de ene eigenschap die niet door één
insluiting gedragen wordt. Lengte nul: $C \subseteq C_n$ met totale
lengte $(\frac23)^n$ (vraag 3). Nergens \omterm{def:b1:topology:dense}{dicht}: $C \subseteq C_n$ dwingt
\omterm{prop:b1:reals:intervals}{intervallen} binnen $C$ tot lengte $\leq 3^{-n}$ (vraag 4). De
overaftelbaarheid rijdt op geen enkele insluiting: ze heeft de volledige
doorsnedestructuur nodig, gecodeerd in de bijectie van vraag 8.

\textbf{11.} Klap $d_n$ om tot $2 - d_n$: de nieuwe rij is nog steeds een
rij over $\{0,2\}$, dus ligt haar waarde $x_n$ in $C$; de partiële sommen
voorbij rang $n$ verschillen precies $2\cdot3^{-n}$, dus $\abs{x_n - x} =
2\cdot3^{-n}$. Bijgevolg is $x_n \neq x$ en $x_n \to x$: elk punt van $C$
is een limiet van andere punten van $C$ --- $C$ is \emph{perfect}, zonder
geïsoleerd punt.

\textbf{12.} Wegens de inductie van vraag 5 zijn de segmenten van $C_n$
precies de $\intcc{t}{t + 3^{-n}}$, waarbij $t$ de waarden van rijen over
$\{0,2\}$ van lengte $n$ doorloopt. Is $x \in C$ met code $(d_k)$, dan is
de afknotting $t_n$ (de cijfers $d_1 \dots d_n$, dan nullen) dus een
linkereindpunt, met $0 \leq x - t_n \leq 3^{-n}$: de eindpunten liggen
\omterm{def:b1:topology:dense}{dicht} in $C$. Ze vormen een deelverzameling van $\{m/3^n : m, n\}$, een
\omterm{def:b1:logic:sets}{verzameling} geïndexeerd door paren gehele getallen, en dus aftelbaar
(zoals voor $\Q$ in \cref{exo:b1:topology:10}). Omdat $C$ overaftelbaar
is (vraag 9), zijn op aftelbaar veel na alle punten van $C$ geen
eindpunt --- $\frac14$ (vraag 7) is de zichtbare top van die ijsberg.

\textbf{13.} Kies $n$ met $3^{-n} < y - x$. Beide punten $x, y$ liggen in
$C_n$, en ze kunnen niet in hetzelfde segment liggen (lengte $3^{-n} < y
- x$): het verwijderde gat tussen hun segmenten levert een $z$ met $x < z
< y$ en $z \notin C_n \supseteq C$. Elke twee punten van $C$ worden dus
door het complement gescheiden: de enige convexe deelverzamelingen van
$C$ zijn de singletons --- $C$ is totaal onsamenhangend.

\textbf{14.} Is $(d_k)$ de $1$-vrije code van $x$, dan is de rij $(2 -
d_k)$ opnieuw een rij over $\{0,2\}$, met partiële sommen
\[
\sum_{k=1}^{n} (2 - d_k)3^{-k} = (1 - 3^{-n}) - \sum_{k=1}^n
d_k 3^{-k} \longrightarrow 1 - x :
\]
dus $1 - x \in C$. Bijgevolg is $1 - C \subseteq C$, en de \omterm{def:b1:logic:map}{afbeelding}
tweemaal toepassen geeft $1 - C = C$: de \omterm{pb:b1:topology:1}{Cantorverzameling} is symmetrisch
om $\frac12$.

\textbf{15.} De partiële sommen $t_n \to x$ en $t'_n \to x'$ (stijgende
rijen convergeren naar hun \omterm{def:b1:reals:bounds}{supremum}, dus naar de waarde), zodat $t_n +
t'_n \to x + x'$ volgens \cref{thm:b1:seq:operations}. Kies bij gegeven
$y \in \intcc{0}{1}$ met code $(e_k)$ het paar $(a_k, b_k) = (0,0),
(0,2), (2,2)$ naargelang $e_k = 0, 1, 2$: dan is $a_k + b_k = 2e_k$, de
rijen $(a_k)$ en $(b_k)$ zijn rijen over $\{0,2\}$ met waarden $x, x' \in
C$, en
\[
x + x' = \lim_n\,(t_n + t'_n) = \lim_n 2\sum_{k=1}^n e_k 3^{-k}
= 2y :
\]
elke $y \in \intcc{0}{1}$ is het midden van twee punten van $C$.

\textbf{16.} Vraag 15 geeft $\intcc{0}{2} = 2\,\intcc{0}{1} \subseteq C +
C$, en $C + C \subseteq \intcc{0}{1} + \intcc{0}{1} = \intcc{0}{2}$:
gelijkheid. Vervolgens, met $1 - C = C$:
\[
C - C = C + (C - 1) = (C + C) - 1 = \intcc{-1}{1} .
\]
Een concreet geval: $1 = \frac14 + \frac34$, een som van twee leden van
$C$ die geen eindpunt zijn.

\textbf{17.} Lengte meet hoeveel van de rechte de \omterm{def:b1:logic:sets}{verzameling} zelf
inneemt; ze zegt niets over de \omterm{def:b1:logic:sets}{verzameling} \emph{sommen}, die het beeld
is van de familie met twee parameters $C \times C$ onder $(x, x') \mapsto
x + x'$ --- de twee cijferrijen worden onafhankelijk gekozen, en precies
die vrijheid vult $\intcc{0}{2}$. Geen enkele stelling begrenst de lengte
van een somverzameling met de lengtes van de summanden, en $C$ is het
bewijs dat geen enkele dat kan.

\textbf{18.} Volgens \cref{pb:b1:reals:1} (vraag 18) is $x$ rationaal
precies wanneer haar eigenlijke ontwikkeling uiteindelijk periodiek is.
De $1$-vrije code van $x \in C$ is ofwel die eigenlijke ontwikkeling
ofwel de oneigenlijke tweeling van een afbrekende; een afbrekende rij en
haar tweeling (uiteindelijk constant $2$) zijn beide uiteindelijk
periodiek, zodat de periodiciteit van de $1$-vrije code equivalent is met
de rationaliteit van $x$. Staartdeling van $\frac1{13}$ in grondtal $3$
($r_0 = 1$): $3 = 13\cdot0 + 3$, $9 = 13\cdot0 + 9$, $27 = 13\cdot2 + 1$,
en de rest keert terug naar $1$: cijfers $\overline{002}$, dus
$\frac1{13} = (0.\overline{002})_3$, $1$-vrij en periodiek: een rationaal
lid van $C$. (Controle: $\frac{2/27}{1 - 1/27} = \frac{2}{26} =
\frac1{13}$.)

\textbf{19.} De rij met $d_k = 2$ op de driehoeksposities $k =
\frac{j(j+1)}{2}$ en $0$ elders is een rij over $\{0,2\}$, zodat haar
waarde $x^*$ tot $C$ behoort (vraag 8). Ze heeft oneindig veel tweeën met
gaten $j + 1 \to \infty$ tussen opeenvolgende, dus is ze niet
uiteindelijk periodiek (een periode $T$ zou uiteindelijk tweeën met gaten
$\leq T$ afdwingen: het argument van de groeiende gaten uit
\cref{pb:b1:reals:1}, vraag 20); volgens vraag 18 is $x^* \notin \Q$. En
volgens vraag 9 plus de aftelbaarheid van $\Q$ zijn op aftelbaar veel na
alle leden van $C$ irrationaal: $x^*$ is de regel, niet de uitzondering.

\textbf{20.} Zij $y \in \intcc{0}{1}$: ze heeft een binaire code $(c_k)$
met $c_k \in \{0, 1\}$ (\cref{pb:b1:reals:1}, vraag 9, in grondtal $2$;
$y = 1$ neemt de rij van louter enen). Dan is $(2c_k)$ een rij over
$\{0,2\}$, ligt haar waarde $x$ in $C$, en is $h(x)$ de waarde van
$(c_k)$, namelijk $y$: $h$ beeldt $C$ op $\intcc{0}{1}$ af. Was $C$ het
beeld van een \omterm{def:b1:logic:map}{afbeelding} vanuit $\N^*$, dan zou samenstellen met $h$ heel
$\intco{0}{1}$ opsommen, in tegenspraak met de diagonaalstelling uit
\cref{pb:b1:reals:1} (vraag 22): $C$ is opnieuw overaftelbaar. Een
\omterm{def:b1:logic:sets}{verzameling} van lengte nul die \omterm{def:b1:logic:inj}{surjectief} is op een heel segment.

\textbf{21.} Volgens vraag 5 is $C$ de disjuncte vereniging van
$\varphi_0(C)$ en $\varphi_2(C)$, elk een verschuiving van de geschaalde
kopie $\frac13 C$. Additiviteit, schaling en translatie-invariantie geven
\[
L(C) = L(\varphi_0(C)) + L(\varphi_2(C))
= \tfrac13 L(C) + \tfrac13 L(C) = \tfrac23\,L(C),
\]
dus $\frac13 L(C) = 0$: $L(C) = 0$. De zelfgelijkvormigheid alleen al
veroordeelt $C$ tot lengte nul --- vraag 3 voltrok slechts het vonnis.

\textbf{22.} Een segment van lengte $l_n$ verliest een centraal \omterm{prop:b1:reals:intervals}{interval}
van lengte $4^{-(n+1)}$ en laat twee segmenten van lengte $l_{n+1} =
\frac{l_n - 4^{-(n+1)}}{2}$ over; vanaf $l_0 = 1$ bevestigt inductie $l_n
= \frac{2^n + 1}{2\cdot4^n}$: immers
$\frac12\Bigl(\frac{2^n+1}{2\cdot4^n} - \frac{1}{4^{n+1}}\Bigr) =
\frac{2(2^n + 1) - 1}{2\cdot4^{n+1}} = \frac{2^{n+1} +
1}{2\cdot4^{n+1}}$, en steeds is $l_n > 0$: de constructie verhongert
nooit. $K = \bigcap K_n$ is \omterm{def:b1:topology:closed}{gesloten} en begrensd en dus compact; een
\omterm{prop:b1:reals:intervals}{interval} binnen $K$ ligt in één segment van $K_n$, van lengte $l_n \to
0$: leeg \omterm{def:b1:topology:closure}{inwendige}. Verwijderde lengte: $\sum_{n\geq0} 2^n \cdot
4^{-(n+1)} = \frac14\sum_{n\geq0}\bigl(\frac12\bigr)^n = \frac12$, en
elke $K_n$ heeft totale lengte $2^n l_n = \frac{2^n + 1}{2^{n+1}} >
\frac12$. Zij nu de vereniging $U$ van eindig veel \omterm{def:b1:topology:open}{open} \omterm{prop:b1:reals:intervals}{intervallen} met
$U \supseteq K$. Volgens de metgezeluitspraak uit
\cref{ex:b1:topology:nested} is $U \supseteq K_n$ voor een zekere $n$;
nemen we de additiviteit van de lengte op eindige verenigingen van
\omterm{prop:b1:reals:intervals}{intervallen} aan, dan is de totale lengte van de overdekkende \omterm{prop:b1:reals:intervals}{intervallen}
minstens die van $K_n$, die $\frac12$ overtreft. $K$ is dus nergens \omterm{def:b1:topology:dense}{dicht}
en toch bestaat er geen goedkope overdekking: topologische kleinheid
(nergens \omterm{def:b1:topology:dense}{dicht}) en metrische kleinheid (lengte nul) zijn werkelijk
verschillende begrippen, en $K$ scheidt ze.

\textbf{23.} Bereiken: zij $d = d(x, F)$ en kies $a_k \in F$ met $\abs{x
- a_k} \leq d + \frac1k$: de $a_k$ zijn begrensd, dus extraheert
Bolzano--Weierstrass (\cref{thm:b1:seq:bw}) een $a_{\varphi(k)} \to a$,
met $a \in F$ ($F$ \omterm{def:b1:topology:closed}{gesloten}, \cref{thm:b1:topology:seqclosed}) en $\abs{x
- a} = \lim \abs{x - a_{\varphi(k)}} = d$. Nu het maximum: is $y \in C$,
dan is $d(y, C) = 0$; anders ligt $y$ in een gat dat in een zeker stadium
$n \geq 1$ verwijderd werd, een \omterm{def:b1:topology:open}{open} \omterm{prop:b1:reals:intervals}{interval} van lengte $3^{-n}$
waarvan de twee eindpunten tot $C$ behoren (vraag 2), dus $d(y, C) \leq
\frac{3^{-n}}{2} \leq \frac16$, met gelijkheid alleen wanneer $n = 1$ en
$y$ het middelpunt van het gat $\intoo{\frac13}{\frac23}$ is, dus $y =
\frac12$; en inderdaad is $d\bigl(\frac12, C\bigr) = \frac16$ omdat $C
\cap \intoo{\frac13}{\frac23} = \emptyset$ en $\frac13, \frac23 \in C$.
Bijgevolg is $\max_{y\in\intcc{0}{1}} d(y, C) = \frac16$, precies bereikt
in $\frac12$.

\textbf{24.} De eindpunten vormen een aftelbare \omterm{def:b1:topology:dense}{dichte} deelverzameling
van $C$ (vraag 12): som ze op als één rij $(e_j)_{j\geq1}$, een rij in
$C$. Al haar deelrijlimieten liggen in $C$ ($C$ \omterm{def:b1:topology:closed}{gesloten}). Omgekeerd, leg
$x \in C$ vast: voor elke $n$ hebben de segmenten van de $C_m$ die $x$
bevatten ($m \geq n$) hun eindpunten binnen $3^{-m} \leq 3^{-n}$ van $x$,
dus liggen er oneindig veel verschillende eindpunten binnen $3^{-n}$ van
$x$; kies indices $j_1 < j_2 < \dots$ met $\abs{e_{j_n} - x} \leq
3^{-n}$: een \omterm{def:b1:seq:subsequence}{deelrij} die naar $x$ convergeert. De \omterm{def:b1:logic:sets}{verzameling}
deelrijlimieten van $(e_j)$ is dus precies $C$ --- één rij die zich bij
overaftelbaar veel punten ophoopt, het tegenovergestelde uiterste van een
convergente rij, waarvan de klonterverzameling een singleton is.

\textbf{25.} (i) De constructie verbruikte: de stabiliteit van \omterm{def:b1:topology:closed}{gesloten
verzamelingen} onder willekeurige doorsnede (het bestaan van $C$ als
\omterm{def:b1:topology:closed}{gesloten verzameling}), de compactheidsstelling
\cref{thm:b1:topology:compact} (vragen 2, 22 en 23), en de
karakteriseringen met rijen van geslotenheid en adherentie (het argument
met geneste compacte \omterm{def:b1:logic:sets}{verzamelingen} uit \cref{ex:b1:topology:nested} en
vraag 23). (ii) De vier paren: lengte nul en toch overaftelbaar (vragen 3
en 9); \omterm{def:b1:topology:closed}{gesloten} en toch met leeg \omterm{def:b1:topology:closure}{inwendige} (vraag 4); perfect --- geen
enkel geïsoleerd punt --- en toch totaal onsamenhangend (vragen 11 en
13); verwaarloosbaar en toch $C + C = \intcc{0}{2}$ (vraag 16). (iii) De
surjectie $h$ uit vraag 20 wordt, continu en niet-dalend gemaakt, de
duivelstrap in de theorie van de continue functies; en in de maattheorie
van het volume van bachelorjaar 3 is $C$ de standaardgetuige dat
``verwaarloosbaar'' niet ``aftelbaar'' betekent, met haar dikke neef
(vraag 22) die ``nergens \omterm{def:b1:topology:dense}{dicht}'' van ``verwaarloosbaar'' scheidt.
\end{solution}
